% Options for packages loaded elsewhere
\PassOptionsToPackage{unicode}{hyperref}
\PassOptionsToPackage{hyphens}{url}
\PassOptionsToPackage{dvipsnames,svgnames,x11names}{xcolor}
%
\documentclass[
  a4paper,
  13pt]{extreport}
\usepackage{amsmath,amssymb}
\usepackage{setspace}
\usepackage{iftex}
\ifPDFTeX
  \usepackage[T1]{fontenc}
  \usepackage[utf8]{inputenc}
  \usepackage{textcomp} % provide euro and other symbols
\else % if luatex or xetex
  \usepackage{unicode-math} % this also loads fontspec
  \defaultfontfeatures{Scale=MatchLowercase}
  \defaultfontfeatures[\rmfamily]{Ligatures=TeX,Scale=1}
\fi
\usepackage{lmodern}
\ifPDFTeX\else
  % xetex/luatex font selection
  \setmainfont[]{Times New Roman}
\fi
% Use upquote if available, for straight quotes in verbatim environments
\IfFileExists{upquote.sty}{\usepackage{upquote}}{}
\IfFileExists{microtype.sty}{% use microtype if available
  \usepackage[]{microtype}
  \UseMicrotypeSet[protrusion]{basicmath} % disable protrusion for tt fonts
}{}
\usepackage{xcolor}
\usepackage[left=3cm,right=2cm,top=2cm,bottom=2cm,headsep=1cm,footskip=1cm]{geometry}
\usepackage{longtable,booktabs,array}
\usepackage{calc} % for calculating minipage widths
% Correct order of tables after \paragraph or \subparagraph
\usepackage{etoolbox}
\makeatletter
\patchcmd\longtable{\par}{\if@noskipsec\mbox{}\fi\par}{}{}
\makeatother
% Allow footnotes in longtable head/foot
\IfFileExists{footnotehyper.sty}{\usepackage{footnotehyper}}{\usepackage{footnote}}
\makesavenoteenv{longtable}
\setlength{\emergencystretch}{3em} % prevent overfull lines
\providecommand{\tightlist}{%
  \setlength{\itemsep}{0pt}\setlength{\parskip}{0pt}}
\setcounter{secnumdepth}{-\maxdimen} % remove section numbering
% ctu_xelatex_header.tex — header-includes cho build LaTeX CTU (xelatex).
% Ánh xạ các ký tự Unicode mà TeX Gyre Termes thiếu glyph (chỉ số dưới/trên,
% ký tự kẻ hộp) sang lệnh LaTeX tương đương để không bị rớt glyph khi biên dịch.
\usepackage{newunicodechar}
\usepackage{pifont}
\usepackage{amssymb}
\newunicodechar{✓}{\ding{51}}
\newunicodechar{✗}{\ding{55}}
\newunicodechar{≥}{\ensuremath{\geq}}
\newunicodechar{≤}{\ensuremath{\leq}}
\newunicodechar{≈}{\ensuremath{\approx}}
\newunicodechar{↔}{\ensuremath{\leftrightarrow}}
\newunicodechar{₀}{\textsubscript{0}}
\newunicodechar{₁}{\textsubscript{1}}
\newunicodechar{₂}{\textsubscript{2}}
\newunicodechar{₃}{\textsubscript{3}}
\newunicodechar{₄}{\textsubscript{4}}
\newunicodechar{₅}{\textsubscript{5}}
\newunicodechar{₆}{\textsubscript{6}}
\newunicodechar{₇}{\textsubscript{7}}
\newunicodechar{₈}{\textsubscript{8}}
\newunicodechar{₉}{\textsubscript{9}}
\newunicodechar{²}{\textsuperscript{2}}
\newunicodechar{³}{\textsuperscript{3}}
\newunicodechar{─}{\rule[0.5ex]{0.9em}{0.4pt}}

% --- Bảng rộng: thu nhỏ cỡ chữ + khoảng cột, cho phép ngắt token dài để
%     bảng không tràn lề phải (overfull hbox). Áp dụng cho mọi longtable
%     (pandoc sinh toàn bộ bảng dạng longtable).
\usepackage{etoolbox}
\usepackage{seqsplit}      % ngắt được token rất dài (mã biến, số dài)
% Lưu ý: KHÔNG đổi \tabcolsep — pandoc tính bề rộng cột theo công thức
% (\columnwidth - 4\tabcolsep)*frac (tham chiếu \tabcolsep trực tiếp), nên
% giảm \tabcolsep lại làm cột RỘNG hơn và tràn nặng hơn.
\AtBeginEnvironment{longtable}{%
  \footnotesize
  \emergencystretch=3em
  \hbadness=10000          % không cảnh báo underfull do RaggedRight
}
% Cho phép TeX bẻ dòng linh hoạt hơn trong ô p{} hẹp
\setlength{\emergencystretch}{2em}

% URL/DOI dài: cho phép ngắt ở bất kỳ ký tự nào (tránh tràn lề ở danh mục TLTK).
% Kết hợp với pandoc +autolink_bare_uris để URL trần được bọc \url{}.
\usepackage{xurl}
\ifLuaTeX
  \usepackage{selnolig}  % disable illegal ligatures
\fi
\IfFileExists{bookmark.sty}{\usepackage{bookmark}}{\usepackage{hyperref}}
\IfFileExists{xurl.sty}{\usepackage{xurl}}{} % add URL line breaks if available
\urlstyle{same}
\hypersetup{
  colorlinks=true,
  linkcolor={black},
  filecolor={Maroon},
  citecolor={black},
  urlcolor={blue},
  pdfcreator={LaTeX via pandoc}}

\author{}
\date{}

\begin{document}

\setstretch{1.2}
\hypertarget{internationalization-and-firm-performance-in-asia-an-institution-contingent-reconception}{%
\section{Internationalization and Firm Performance in Asia: An
Institution-Contingent
Reconception}\label{internationalization-and-firm-performance-in-asia-an-institution-contingent-reconception}}

\hypertarget{integrative-synthesis-essay-kappa}{%
\subsection{Integrative Synthesis Essay
(Kappa)}\label{integrative-synthesis-essay-kappa}}

\emph{Đỗ Thùy Hương, Doctoral dissertation, Can Tho University}
\emph{Supervisor: Assoc. Prof.~Dr.~Phan Anh Tú}

\begin{quote}
This essay is the English-language \textbf{integrative synthesis
(kappa)} of a Vietnamese-language doctoral dissertation. Following the
compilation-thesis convention, it does not repeat the component studies
in full; it states the single argument that binds them, shows how each
component earns its place in that argument, and articulates the
theoretical contribution at a register intended for an international
international-business (IB) readership. All quantitative results are
reported as they appear in the dissertation's results chapters and are
not re-derived here.
\end{quote}

\begin{center}\rule{0.5\linewidth}{0.5pt}\end{center}

\hypertarget{phenomenon-and-motivation}{%
\subsection{1. Phenomenon and
motivation}\label{phenomenon-and-motivation}}

For four decades the internationalization--performance (I--P)
relationship has been one of the most heavily researched (and most
stubbornly inconsistent) questions in international business.
Meta-analyses recover effect sizes ranging from strongly positive to
null to negative, and functional-form claims span linear, U-shaped,
inverted-U, S-shaped, and horizontal-J specifications. The dominant
reconciliation, since Lu and Beamish (2004) and Contractor, Kundu and
Hsu (2003), has been the \textbf{inverted-U}: early internationalization
incurs liabilities of foreignness and coordination costs, an
intermediate range yields scale, learning, and market-diversification
benefits, and beyond a turning point governance costs overwhelm those
benefits. Yet the inverted-U is itself unstable across contexts, and the
field has largely treated this instability as residual noise (``context
matters'') rather than as a structured, predictable phenomenon.

Three contemporary forces sharpen the puzzle and motivate this
dissertation. First, an \textbf{employment and productivity imperative}
across developing Asia makes the firm-level returns to
internationalization a first-order policy question, not merely an
academic one. Second, an \textbf{institutional spectrum} of unusual
breadth coexists within a single region, ranging from innovation-driven
advanced economies (Singapore, Korea, Japan) through resource-rich Gulf
states, large emerging markets (China), lower-middle-income transition
economies (Vietnam), to frontier economies and Pacific Small Island
Developing States (SIDS). Third, the \textbf{digital era} has altered
the learning mechanism at the heart of internationalization theory,
raising the question of whether digital capability changes the shape,
not merely the level, of the I--P relationship. These three forces
converge on a firm-level question that the existing literature answers
only fragmentarily.

\hypertarget{research-gap-and-questions}{%
\subsection{2. Research gap and
questions}\label{research-gap-and-questions}}

The gap this dissertation addresses is not the absence of evidence but
the absence of \textbf{structure} in that evidence. Prior work has shown
that institutions moderate the \emph{magnitude} of the I--P effect
(Marano et al., 2016); it has not shown whether institutions also shape
the \emph{sign and functional form} of the effect across a continuous
spectrum, nor located the regime at which the inverted-U family may
cease to apply. In parallel, the digitalisation literature has tended to
conflate distinct constructs (deep technological capability versus
surface-level digital adoption), obscuring their different mechanisms.

The dissertation therefore poses a central question, \emph{How does
internationalization affect the performance of Asian firms, and what
institutional, digital, and managerial conditions explain the high
heterogeneity in that relationship?}, decomposed into four research
questions:

\begin{itemize}
\tightlist
\item
  \textbf{RQ1.} What is the pooled effect of internationalization on
  firm performance in the 1977--2026 literature, how heterogeneous is
  it, and which moderators systematically explain that heterogeneity?
\item
  \textbf{RQ2.} Within specific national contexts (Singapore, Vietnam,
  China), through which mechanisms does the I--P relationship operate
  and how does its nonlinearity manifest?
\item
  \textbf{RQ3.} Across 50 Asia-Pacific economies (including seven
  Pacific SIDS), do the moderating roles of digital adoption (DAI) and
  technological capability (TCI) generalise, and do these two constructs
  in fact operate through different mechanisms?
\item
  \textbf{RQ4.} What is the role of institutions (the ICRV sub-regimes),
  of the TCI--DAI distinction, and of top-manager characteristics in
  moderating the I--P relationship, and what does their three-way
  interaction imply?
\end{itemize}

\hypertarget{theoretical-framework-the-cdcm-architecture}{%
\subsection{3. Theoretical framework: the CDCM
architecture}\label{theoretical-framework-the-cdcm-architecture}}

The dissertation's theoretical contribution is organised by a single
integrative architecture, the \textbf{Country--Digital--Capability
Moderation (CDCM)} framework, which combines four classic theories with
a digital-capability lens not as a mechanical concatenation but as an
internally consistent multi-tier explanation in which each tier accounts
for a distinct source of heterogeneity in the I--P relationship.

\begin{itemize}
\tightlist
\item
  The \textbf{Uppsala internationalization model} supplies the baseline
  learning logic and, extended for the AI era, a ``data-augmented
  learning'' mechanism (Banalieva \& Dhanaraj, 2019).
\item
  The \textbf{resource-based view} (Wernerfelt, 1984; Barney, 1991)
  frames technological capability as an absorptive, depth-of-capability
  resource that widens the adaptation frontier.
\item
  \textbf{Institutional theory} (North, 1990; Khanna \& Palepu, 2010)
  supplies the macro tier and, crucially, the claim that institutional
  context conditions firm strategy outcomes.
\item
  \textbf{Upper-echelons theory} (Hambrick \& Mason, 1984) supplies the
  individual tier of manager experience and diversity.
\end{itemize}

The macro tier is operationalised through the \textbf{Institutional
Context Regime Variation (ICRV)} classification, a six-regime taxonomy
(I innovation-driven advanced; II resource-based advanced; III
upper-middle; IV lower-middle transition; V emerging; VI SIDS) that
converts a continuous institutional spectrum into a testable firm-level
moderator. The contribution of ICRV is not the addition of a categorical
control but the provision of a systematic way to predict how the turning
point of the I--P curve shifts with institutional regime, turning the
qualitative claim ``context matters'' into a falsifiable quantitative
proposition. The architecture's limiting case is the \textbf{Forced
Internationalization Penalty (FIP)}, a named boundary condition
proposing that when all three CDCM tiers are simultaneously minimal, the
I--P relationship may lose not only its early-stage benefit but its
inverted-U form, tilting toward a monotone-negative limiting case (the
empirical magnitude of this tilt is data-version-sensitive and treated
as exploratory; see the SIDS evidence below).

\hypertarget{the-compilation-and-how-each-component-earns-its-place}{%
\subsection{4. The compilation and how each component earns its
place}\label{the-compilation-and-how-each-component-earns-its-place}}

The dissertation is a compilation of ten component studies plus a
supplementary book chapter, spanning a meta-analysis, five
single-country estimations distributed across the institutional
spectrum, one 50-economy pooled test, and two boundary cases. The
empirical backbone is harmonised World Bank Enterprise Surveys (WBES)
microdata across three questionnaire generations (PICS3, Standardized,
BREADY/BEE), 2006--2026. The frame hierarchy is explicit: a
classification pool of 96,415 firms / 52 labels narrows to an analytic
frame of 88,869 firms / 50 economies / 103 economy-year cells, within
which the labour-productivity-valid descriptive frame is 84,998 firms /
107 cells, and the capstone P7 regression frame is 81,022 (M2) / 79,080
(M5).

Table 1 traces the \textbf{golden thread}: the line that runs from each
research question, through the components that answer it and the
hypotheses they test, to the empirical finding and the theoretical
contribution it supports. Read across, each row shows the internal
consistency of one link; read down, the rows show how independent pieces
of evidence accumulate into one claim.

\textbf{Table 1.} \emph{Golden thread: research question → component
(ICRV regime) → hypothesis → key finding → theoretical contribution.}

\begin{longtable}[]{@{}
  >{\raggedright\arraybackslash}p{(\columnwidth - 8\tabcolsep) * \real{0.2000}}
  >{\raggedright\arraybackslash}p{(\columnwidth - 8\tabcolsep) * \real{0.2000}}
  >{\raggedright\arraybackslash}p{(\columnwidth - 8\tabcolsep) * \real{0.2000}}
  >{\raggedright\arraybackslash}p{(\columnwidth - 8\tabcolsep) * \real{0.2000}}
  >{\raggedright\arraybackslash}p{(\columnwidth - 8\tabcolsep) * \real{0.2000}}@{}}
\toprule\noalign{}
\begin{minipage}[b]{\linewidth}\raggedright
RQ
\end{minipage} & \begin{minipage}[b]{\linewidth}\raggedright
Component (ICRV)
\end{minipage} & \begin{minipage}[b]{\linewidth}\raggedright
Hypothesis
\end{minipage} & \begin{minipage}[b]{\linewidth}\raggedright
Key empirical finding
\end{minipage} & \begin{minipage}[b]{\linewidth}\raggedright
Theoretical contribution
\end{minipage} \\
\midrule\noalign{}
\endhead
\bottomrule\noalign{}
\endlastfoot
Foundation & P1 --- 17-economy Asian baseline (\emph{N} = 40,633) &
Foundational (CDCM) & Regional productivity stratification;
\textbf{digital-shield effect} (obstacle × technology +0.110, \emph{p}
\textless{} .001) & Seeds CDCM; origin of the ``digital shield'' that P7
scales to 50 economies \\
RQ1 & P6 --- global meta-analysis & H1, H5 & \emph{k} = 238; small
positive pooled \emph{r̄} = 0.074; high heterogeneity \emph{I²} = 62.4\%;
institutions moderate \emph{Q}M = 17.35 (\emph{p} = .002) & Establishes
that I--P is small-positive but \textbf{institution-driven
heterogeneous} --- empirical basis for the ICRV axis \\
RQ2 & P4 --- Singapore (I) & H1, H3 & Near-linear positive; turning
point \textasciitilde88.6\% (wide CI, outside range); DAI amplifies at
high intensity & Upper bound: the curve \textbf{straightens} when the
institutional cushion is thick \\
RQ2 & P3 --- Vietnam (IV) & H1, H2, H3 & Sharp inverted-U, turning point
39.3--46.2\%; TCI positively moderates & Anchor regime; distinguishes
the \textbf{participation margin} from the intensity margin \\
RQ2 & P5 --- China 2012--2024 (III); P2 --- Chinese SMEs & H1, H6 &
Turning point stable \textasciitilde48.78\% (Paternoster \emph{z} =
0.82, \emph{p} = .412); P2: 47.8\%, working-capital trap on the
downslope & \textbf{Temporal robustness} of the turning point across a
decade \\
RQ3 & P7 --- 50-economy pooled & H1, H2, H3, H5 & \emph{N} = 81,022 /
79,080; turning point 43.6--51.5\%; TCI +0.141, DAI +0.201 (\emph{p}
\textless{} .001) raise the level & Generalises CDCM; TCI/DAI act as
\textbf{universal level-shifters}, not curve-benders \\
RQ3 & P8 --- Pacific SIDS (VI) & H1b & Structure loss: slope −0.085
(\emph{p}wild = .66), curvature +0.696 (\emph{p}wild = .082, marginal);
\emph{N} = 1,450 & Names \textbf{FIP} --- the extreme boundary condition
of the inverted-U \\
RQ4 & P7 --- full ICRV spectrum & H5 & Turning point shifts by regime
(Group IV 43.0\% → Group I near-linear → Groups V/VI structure loss) &
Institutions moderate the \textbf{functional form}, not only the
magnitude --- ``context-contingent'' \\
RQ4 & P9 --- India 2014--2025 (IV) & H1c & Turning point
\textbf{dissolves}: 61.8\% (2014) → 40.7\% (2022) → curvature
insignificant (2025: \emph{β₂} = −0.16, \emph{p} = .42) & Intense
institutional transition can \textbf{dissolve}, not merely shift, the
threshold (temporal boundary) \\
RQ4 & P10 --- Japan 2025 (I) & H1 (upper bound) & Near-linear positive,
turning point outside range; export premium 25.8\% & Confirms the
\textbf{upper-bound prediction}; participation-margin mechanism (Melitz,
2003) \\
RQ4 & Book chapter --- India, governance (IV) & H4 & ROS: experience
+1.55, female leadership +77.03 positively moderate (\emph{N} = 380) &
Single-country evidence for the \textbf{managerial tier} of CDCM \\
\end{longtable}

\emph{Note. Figures are reproduced from the dissertation's results
chapters (§4.2--§4.9; §5.6.1). ICRV groups: I = innovation-driven
advanced; III = upper-middle; IV = lower-middle transition; VI = SIDS.}

\hypertarget{the-integrated-finding-a-three-zone-morphology}{%
\subsection{5. The integrated finding: a three-zone
morphology}\label{the-integrated-finding-a-three-zone-morphology}}

Placed side by side, the full set of ten component studies plus the
supplementary book chapter converge on a single, falsifiable theoretical
claim. This set comprises the 17-economy foundational study (P1), the
global meta-analysis (P6), five single-country estimations (P2--P5, P9),
the 50-economy pooled test (P7), two boundary cases (P8 SIDS, P10
Japan), and the governance-tier book chapter. The claim is this:
\textbf{internationalization improves firm performance through a
functional form that is itself a function of institutional position.}
The thickness of the institutional cushion determines whether the I--P
relationship is near-linear positive, a sharp inverted-U, or
structurally absent and sign-reversed.

Concretely, a \textbf{three-zone morphology} organises the entire body
of evidence as a continuous expression of one institutional mechanism
rather than three separate phenomena:

\begin{itemize}
\tightlist
\item
  \textbf{Upper bound} (Group I, innovation-driven advanced; Singapore
  P4, Japan P10): the downslope is pushed out beyond the feasible export
  range, so within the data the relationship is near-linear positive;
  the mechanism shifts from the intensity margin to the participation
  margin (export-versus-not), consistent with Melitz-type
  self-selection.
\item
  \textbf{Anchor regimes} (Groups III--IV, upper-middle and lower-middle
  transition; China P5/P2, Vietnam P3, India P9): the inverted-U is
  sharpest and best-identified, with turning points clustering around
  40--49\%; independent single-country estimates converge on this band
  without being constrained to do so, an out-of-sample coincidence that
  resists a specification-search reading.
\item
  \textbf{Lower bound} (Groups V--VI, emerging and SIDS; P7, P8):
  structure is lost. In Group V the coefficients lose significance; in
  the SIDS regime the inverted-U dissolves and the relationship tilts
  toward the monotone-negative FIP limiting case.
\end{itemize}

This evidence points from the field's classic claim that ``context
matters'' (institutions moderate the \emph{magnitude} of the effect)
toward \textbf{context-contingency} (institutions may also moderate the
\emph{sign and form} of the effect), and it locates the regime (the SIDS
extreme) at which the inverted-U form appears to dissolve.

\hypertarget{theoretical-contributions}{%
\subsection{6. Theoretical
contributions}\label{theoretical-contributions}}

Following the criteria of originality × utility (Corley \& Gioia, 2011)
and what/how/why (Whetten, 1989), the dissertation makes four
contributions, each stated as a falsifiable proposition with explicit
scope conditions.

\begin{enumerate}
\def\labelenumi{\arabic{enumi}.}
\item
  \textbf{DAI as a situational resource, distinct from TCI as a
  foundational capability.} The 50-economy evidence shows DAI raising
  the performance level universally (+0.201, \emph{p} \textless{} .001)
  but bending the I--P curve only context-dependently (clearest in the
  strong-institution single-country case, P4), with its moderating
  effect \emph{stronger} in weaker-institution regimes, a ``digital
  shield'' that compensates for institutional deficits. This separates
  two constructs the literature routinely conflates and specifies
  \emph{how} each operates (level-shift versus curve-bend).
\item
  \textbf{The Forced Internationalization Penalty (FIP).} A newly named
  boundary condition specifying that when domestic-market size, trade
  costs, and institutional support are simultaneously adverse, the I--P
  relationship loses its inverted-U form. Its value is theoretical
  beyond any single estimate: it demarcates the application boundary of
  a whole family of internationalization theories that tacitly assume
  minimum structural conditions are met. The accompanying SIDS evidence
  (P8) is treated transparently as exploratory and
  data-version-sensitive.
\item
  \textbf{The six-regime ICRV taxonomy.} A region-specific institutional
  classification that converts a continuous spectrum into a testable
  moderator capable of predicting turning-point location by regime: a
  methodological-theoretical instrument with reuse value beyond this
  dissertation.
\item
  \textbf{Institutional stability as an implicit premise of the I--P
  literature.} By identifying both the boundary at which the inverted-U
  dissolves (FIP) and the temporal process by which intense
  institutional transition dissolves the threshold (India P9, H1c), the
  dissertation makes explicit a premise the literature had left tacit.
\end{enumerate}

\hypertarget{boundary-conditions-and-identification}{%
\subsection{7. Boundary conditions and
identification}\label{boundary-conditions-and-identification}}

A contribution that meets international standards must state not only
what it claims but the conditions under which the claim holds and the
inferential basis that supports it (Santangelo \& Verbeke, 2022). The
CDCM framework applies at the \textbf{firm level}, with
internationalization measured as \textbf{export intensity} (the
sales-based dimension of the degree of internationalization; Sullivan,
1994); other dimensions (geographic scope, FDI-based commitment depth)
are out-of-scope boundary conditions to be replicated. It is calibrated
on the \textbf{Asia-Pacific institutional spectrum, 2006--2026};
extension to other regional configurations is an open scope test. It
assumes three minimum structural conditions, and the explicit
identification of their joint violation (the SIDS/FIP case) is precisely
how the framework states the boundary of its own application.

On causal identification the dissertation is consistently cautious: the
design is \textbf{repeated cross-section}, so claims are read as
conditional structural relationships rather than point-identified causal
effects. The credibility of the central claim rests not on any single
estimate but on \textbf{triangulation across five identification sources
with non-overlapping biases}: (i) instrumental-variable estimation for
digital adoption, whose near-null instrumented coefficient
\emph{deflates} rather than inflates the DAI--performance association;
(ii) Oster bounds on selection on unobservables; (iii) Paternoster
cross-wave coefficient-stability tests; (iv) sign-consistency across
many specifications as a baseline inferential principle; and (v) the
convergence of the meta-analysis (P6) and the firm-level pooled
estimation (P7), two methods with different bias structures that
nonetheless locate institutions as the pivotal moderator. When five
sources with mutually compensating weaknesses point to one conclusion,
the structural claim becomes hard to attribute to a single analytic
artefact, even where each source alone falls short of point causal
identification.

\hypertarget{implications}{%
\subsection{8. Implications}\label{implications}}

For \textbf{theory}, the dissertation reframes a four-decade debate: the
question is not ``is the I--P relationship positive, negative, or
inverted-U?'' but ``at what institutional position does each form
obtain?'' For \textbf{practice and policy}, it replaces a
one-size-fits-all prescription with a regime-conditioned policy map: in
strong-institution regimes, remove expansion ceilings and raise
absorptive-capacity ceilings; in transition regimes, identify the
optimum and avoid over-expansion ahead of coordination capacity; in the
SIDS regime, address structural trade-cost and institutional constraints
before expecting internationalization to improve performance, because
there internationalization is a survival burden, not a
competitive-advantage choice.

\hypertarget{limitations-and-future-research}{%
\subsection{9. Limitations and future
research}\label{limitations-and-future-research}}

The principal limitations are stated as study weaknesses distinct from
the framework's boundary conditions: the single-dimension
(export-intensity) operationalisation of internationalization; the
labour-productivity (rather than multidimensional) performance measure;
the repeated-cross-section design; and the platform-level
operationalisation of digital adoption. Each maps to a future-research
direction: multidimensional DOI measures, multidimensional performance,
genuine firm-level panels (the forthcoming WBES core panel), richer
Tier-2 digital measures, and replication of the three-zone morphology
beyond Asia-Pacific.

\hypertarget{concluding-statement}{%
\subsection{10. Concluding statement}\label{concluding-statement}}

The dissertation's single new sentence to the field is this: across the
institutional spectrum of Asia and the Pacific, the form of the
internationalization--performance relationship appears to be shaped by
institutions, such that the inverted-U is plausibly one regime-specific
case within a continuum that runs from near-linear reward at the upper
bound to structural penalty at the lower bound; and the dissertation
locates, and names, the regime at which the inverted-U form appears no
longer to apply.

\begin{center}\rule{0.5\linewidth}{0.5pt}\end{center}

\emph{Citation note: all references appear in the dissertation's APA-7
reference list (\texttt{thesis/04\_references\_apa7.md}). This synthesis
introduces no claims or figures beyond those in the Vietnamese
dissertation chapters.}

\end{document}
